\chapter{Series Expansions: Taylor and Fourier}

\begin{definicion}[Sequence]
    A sequence in a set $X$ is a function $\alpha: \mathbb{N} \to X$. We denote $\alpha(n)$ by $x_n$ and write the sequence as:
    \[
        \{x_n\}_{n\in \mathbb{N}}
    \]
\end{definicion}
\begin{definicion}[Series]
    Given a sequence $\{x_n\}_{n\in \mathbb{N}}$, the series associated to it is the sequence of partial sums:
    \[
        \{s_n\}_{n\in \mathbb{N}} \quad \text{where} \quad s_n = \sum_{k=1}^n x_k
    \]

    The series is written as $\sum\limits_{n\geq 1} x_n$. It is said to be convergent if the sequence of partial sums $\{s_n\}$ converges, and in that case, we define:
    \[
        \sum_{n=1}^\infty x_n = \lim_{n\to \infty} s_n
    \] 
\end{definicion}

An specific type of series is the \textbf{power series}, which has the form:
\[
    \sum_{n=0}^\infty a_n (x - c)^n
\]
where $a_n\in \mathbb{R}$ (or $\mathbb{C}$) are the coefficients, $c$ is the center of the series, and $x$ is the variable. The convergence of power series depends on the value of $x$ and can be analyzed using the radius of convergence.
\begin{definicion}[Radius of Convergence]
    The radius of convergence $R$ of a power series $\sum_{n\geq 0} a_n (x - c)^n$ is defined as:
\[
    R = \frac{1}{\limsup_{n\to \infty} \sqrt[n]{|a_n|}}
\]
If the limit superior is zero, we set $R = \infty$, and if it is infinite, we set $R = 0$.
\end{definicion}
\begin{teo}
    Let $\sum_{n\geq 0} a_n (x - c)^n$ be a power series with radius of convergence $R$. Then:
    \begin{itemize}
        \item The series converges for all $x$ such that $\|x - c| < R$.
        \item The series diverges for all $x$ such that $|x - c| > R$.
        \item The behavior of the series at the boundary points where $|x - c| = R$ must be analyzed separately, as it may converge or diverge depending on the specific coefficients $a_n$.
    \end{itemize}
\end{teo}
\begin{prop}
    Given a power series $\sum\limits_{n\geq 0} a_n (x - c)^n$, if $a_n\neq 0$ for all $n\in \mathbb{N}$ and the following limit exists, then the radius of convergence can be computed as:
    \[
    R = \lim_{n\to \infty} \left|\frac{a_n}{a_{n+1}}\right|
    \]
\end{prop}

\section{Taylor Series}

\begin{definicion}[Taylor Series]
    Let $f:\bb{R}\to \bb{R}$ be a function that is infinitely differentiable at a point $c\in \bb{R}$. The Taylor series of $f$ centered at $c\in \bb{R}$ of order $n$ is the power series given by:
\[
    T_{f,c}(x)^n := \sum_{k=0}^n \frac{f^{(k)}(c)}{k!} (x - c)^k
\]
where $f^{(k)}(c)$ denotes the $k$-th derivative of $f$ evaluated at $c$.
\end{definicion}

\begin{teo}[Taylor's Theorem]
    Let $f:\bb{R}\to \bb{R}$ be a function that is $n+1$ times differentiable on an interval containing $c$ and $x$. Then there exists some $\xi$ between $c$ and $x$ such that:
    \[
    f(x) = T_{f,c}(x)^n + \frac{f^{(n+1)}(\xi)}{(n+1)!} (x - c)^{n+1}
    \]
\end{teo}

In other words, the Taylor series approximates the function $f$ around the point $c$, and the error of this approximation is given by the remainder term:
\[
    R_n(x) = \frac{f^{(n+1)}(\xi)}{(n+1)!} (x - c)^{n+1}
\]

Taylor functions are really used to approximate functions around a specific point, and they are particularly useful when the function is difficult to compute directly.
\subsection{Approximating a Pendulum as a Spring-Mass System}

Consider a simple spring-mass system, where a mass $m$ is attached to a spring with spring constant $k$ (which depends on the stiffness of the spring). The equation of motion for this system is given by Hooke's Law:
\[
    F_R = -kx
\]
where $F_R$ is the restoring force exerted by the spring, and $x$ is the displacement of the mass from its equilibrium position. The negative sign indicates that the force is directed opposite to the displacement, trying to restore the mass to its equilibrium position. We will now consider a pendulum, which consists of a mass $m$ attached to a string of length $L$, swinging under the influence of gravity, as described in Figure~\ref{fig:pendulum}.
\begin{figure}
    \centering
    \begin{tikzpicture}[
        force/.style={-Stealth, thick, green!60!black},
    ]

        % --- Configuración ---
        \def\L{3}        
        \def\angleVal{25} % Ángulo phi
        \def\Fmag{2.5}   

        % --- Coordenadas base ---
        \coordinate (O) at (0,0); 
        \coordinate (V) at (0,-2); % Punto auxiliar en el eje vertical (abajo)
        \coordinate (P) at ({-\angleVal-90}:\L); % Posición de la masa

        % --- Soporte ---
        \fill[gray!40] (-1,0.2) rectangle (1,0);
        \draw[thick] (-1,0) -- (1,0);

        % --- Líneas de referencia y péndulo ---
        \draw[dashed, thin, gray] (O) -- (0,-\L-1); 
        \draw[thick] (O) -- node[midway, above, sloped] {$L$} (P);

        % --- Ángulo superior (usando tu sintaxis pic) ---
        \draw pic[
            draw,
            -Stealth,
            angle radius=1.2cm,
            angle eccentricity=1.3,
            pic text={$\varphi$}
        ] {angle = P--O--V};

        % --- Vectores de Fuerza ---
        \coordinate (FG_end) at ($(P) + (0,-\Fmag)$);
        \coordinate (FN_end) at ($(P) + ({-\angleVal-90}:\Fmag*cos{\angleVal})$);
        \coordinate (FT_end) at ($(P) + ({-\angleVal}:\Fmag*sin{\angleVal})$);

        \draw[force] (P) -- (FG_end) node[below] {$\vec{F}_{\text{G}}$};
        \draw[force, dotted] (P) -- (FN_end) node[above left] {$\vec{F}_{\text{G,nor}}$};
        \draw[force, dotted] (P) -- (FT_end) node[below] {$\vec{F}_{\text{G,tan}}$};

        % Proyecciones
        \draw[dotted] (FG_end) -- (FN_end);
        \draw[dotted] (FG_end) -- (FT_end);

        % --- Ángulo inferior en el triángulo de fuerzas ---
        \draw pic[
            draw,
            -Stealth,
            angle radius=0.8cm,
            angle eccentricity=1.4,
            pic text={$\varphi$}
        ] {angle = FN_end--P--FG_end};

        

        \draw[dotted, blue, thick] (P) arc (270-\angleVal:270:\L) node[midway, above, sloped] {\color{blue}$x$};

        % Punto al final del péndulo
        \fill[red] (P) circle (0.15);

    \end{tikzpicture}
    \caption{Pendulum diagram with forces and angles.}
    \label{fig:pendulum}
\end{figure}

The force acting on the mass in the pendulum is the gravitational force, which can be decomposed into two components: one perpendicular to the string (normal component) and one tangential to the arc of motion (tangential component). The tangential component of the gravitational force is responsible for the motion of the pendulum and is given by:
\[
    F_{\text{G,tan}} = - mg \sin(\varphi)
\]
where the negative sign indicates that the force is directed opposite to the direction of motion, and $\varphi$ is the angle between the string and the vertical. Given that $\varphi$ is given in radians, we can easily obtain its value:
\[
    \varphi\cdot L = x \Longrightarrow \varphi = \frac{x}{L}
\]
Substituting this into the expression for the tangential force, we get:
\[
    F_{\text{G,tan}} = - mg \sin\left(\frac{x}{L}\right)
\]

Now, we can use the Taylor series expansion of the sine function around $0$ (which is the equilibrium position of the pendulum) up to the first order term:
\[
    \sin\left(\frac{x}{L}\right) \approx \sin\left(\dfrac{0}{L}\right)(x - 0)^0 + \cos\left(\dfrac{0}{L}\right)\cdot \dfrac{1}{L}(x - 0)^1 = \dfrac{x}{L}
\]
Substituting this approximation back into the expression for the tangential force, we get:
\[
    F_{\text{G,tan}} \approx - mg \cdot \dfrac{x}{L} = - \dfrac{mg}{L} x
\]

This shows that for small angles (small displacements), the pendulum behaves like a spring-mass system with an effective spring constant $k_{\text{eff}} = \dfrac{mg}{L}$.

\section{Fourier Series}